% Preámbulo: paquetes, colores y estilo del informe

\usepackage[utf8]{inputenc}
\usepackage[T1]{fontenc}
\usepackage[spanish,es-noshorthands,es-nodecimaldot,es-tabla]{babel}
\usepackage[scaled=0.95]{helvet}
\renewcommand{\familydefault}{\sfdefault}
\usepackage{amsmath}
\usepackage{microtype}
\usepackage[strings,nohyphen]{underscore}

\usepackage[a4paper,top=24mm,bottom=22mm,left=18mm,right=18mm,headheight=14pt]{geometry}
\usepackage[table]{xcolor}
\usepackage{graphicx}
\usepackage{tikz}
\usepackage{xltabular}
\usepackage{array}
\usepackage{enumitem}
\usepackage{titlesec}
\usepackage{tocloft}
\usepackage{fancyhdr}
\usepackage{caption}
\usepackage{placeins}
\usepackage{parskip}

% Paleta
\definecolor{azul}{HTML}{153A5B}
\definecolor{azuldos}{HTML}{1E4D73}
\definecolor{texto}{HTML}{1C2834}
\definecolor{gris}{HTML}{5C6B7A}
\definecolor{grisclaro}{HTML}{7B8794}
\definecolor{dorado}{HTML}{C4A15A}
\definecolor{bandaclara}{HTML}{D5E4F2}
\definecolor{fila}{HTML}{F4F7FB}
\definecolor{borde}{HTML}{D5DEE6}
\definecolor{fondofig}{HTML}{FBFCFE}
\definecolor{fondoleyenda}{HTML}{F7F9FB}
\color{texto}

% Párrafos y listas
\setlength{\parskip}{6pt}
\setlist{topsep=2pt,itemsep=2pt,parsep=0pt}

% Títulos
\titleformat{\section}{\color{azul}\Large\bfseries}{\thesection.}{0.5em}{}[{\color{azul}\titlerule[1.2pt]}]
\titleformat{\subsection}{\color{azuldos}\large\bfseries}{\thesubsection}{0.6em}{}
\titlespacing*{\section}{0pt}{16pt}{8pt}
\titlespacing*{\subsection}{0pt}{12pt}{4pt}

% Índice
\addto\captionsspanish{\renewcommand{\contentsname}{Contenido}}
\renewcommand{\cftsecleader}{\cftdotfill{\cftdotsep}}
\renewcommand{\cftsecfont}{\bfseries\color{azul}}
\renewcommand{\cftsecpagefont}{\bfseries\color{azul}}
\renewcommand{\cftsecaftersnum}{.}
\renewcommand{\cftsubsecpagefont}{\bfseries\color{azul}}
\renewcommand{\cfttoctitlefont}{\color{azul}\Large\bfseries}
\renewcommand{\cftaftertoctitle}{\par\vspace{-6pt}{\color{azul}\rule{\linewidth}{1.2pt}}}

% Encabezado y pie
\pagestyle{fancy}
\fancyhf{}
\fancyhead[L]{\footnotesize\color{gris}Diseño de la base de datos}
\fancyhead[R]{\footnotesize\color{gris}MERXIRVE}
\fancyfoot[L]{\scriptsize\color{grisclaro}Bases de Datos II \textperiodcentered{} Ingeniería de Software \textperiodcentered{} 2026-2027}
\fancyfoot[R]{\small\bfseries\color{azul}\thepage}
\renewcommand{\headrulewidth}{0pt}

% Figuras
\captionsetup{font={small,bf},labelsep=period,skip=4pt}
\newcommand{\figura}[3]{%
  \begin{figure}[htbp]
    \centering
    \setlength{\fboxsep}{2pt}%
    \fcolorbox{borde}{fondofig}{\includegraphics[width=\dimexpr\linewidth-2\fboxsep-2\fboxrule\relax,height=0.42\textheight,keepaspectratio]{assets/export/#1}}
    \caption{#2}\label{#3}
  \end{figure}}

% Tablas: cabecera azul y filas alternas
\renewcommand{\arraystretch}{1.3}
\newcolumntype{L}[1]{>{\raggedright\arraybackslash\hsize=#1\hsize}X}
\newcommand{\cab}[1]{\textcolor{white}{\bfseries #1}}
\newcommand{\filas}{\rowcolors{2}{fila}{white}}

% Muestras de la leyenda de notación
\newcommand{\muestra}[1]{\tikz[baseline=-0.6ex]{#1}}
\newcommand{\swEnt}{\muestra{\draw[fill=blue!10,draw=azul,line width=0.9pt] (0,-0.13) rectangle (0.45,0.13);}}
\newcommand{\swRel}{\muestra{\draw[fill=yellow!20,draw=brown!80!black,line width=0.7pt] (0,0) -- (0.16,0.16) -- (0.32,0) -- (0.16,-0.16) -- cycle;}}
\newcommand{\swAgg}{\muestra{\draw[fill=fondoleyenda,draw=gris,line width=0.7pt] (0,-0.17) rectangle (0.62,0.17); \draw[fill=blue!10,draw=azul,line width=0.6pt] (0.06,-0.08) rectangle (0.24,0.08); \draw[fill=yellow!20,draw=brown!80!black,line width=0.5pt] (0.30,0) -- (0.38,0.08) -- (0.46,0) -- (0.38,-0.08) -- cycle;}}
\newcommand{\swAtr}{\muestra{\draw[fill=green!8,draw=green!50!black,line width=0.7pt] (0.22,0) ellipse (0.22 and 0.13);}}
\newcommand{\swClave}{\muestra{\draw[fill=green!8,draw=green!50!black,line width=0.7pt] (0.22,0) ellipse (0.22 and 0.13); \draw[line width=0.5pt] (0.1,-0.05) -- (0.34,-0.05);}}
\newcommand{\swHer}{\muestra{\draw[fill=white,draw=green!50!black,dashed,line width=0.7pt] (0.22,0) ellipse (0.22 and 0.13); \draw[line width=0.5pt] (0.1,-0.05) -- (0.34,-0.05);}}
\newcommand{\simbolo}[2]{\mbox{#1\enspace#2}\hspace{1.2em plus 0.5em}}

\usepackage[hidelinks,bookmarksnumbered]{hyperref}
\hypersetup{pdftitle={Informe: Diseño de la base de datos},pdfauthor={Equipo MERXIRVE}}
